% Compile with XeLaTeX twice. All mathematics uses equation/align environments.
\documentclass[12pt]{article}
\usepackage[margin=1in]{geometry}
\usepackage{amsmath,amssymb,amsthm}
\usepackage{microtype}
\usepackage[hidelinks]{hyperref}
% Load the template's New Computer Modern family directly to avoid
% font-setup warnings in the installed newcomputermodern package.
\usepackage{fontspec}
\setmainfont{NewCM10-Book.otf}[BoldFont=NewCM10-Bold.otf,
 ItalicFont=NewCM10-BookItalic.otf,BoldItalicFont=NewCM10-BoldItalic.otf]
\usepackage{unicode-math}
\setmathfont{NewCMMath-Book.otf}
\usepackage[most]{tcolorbox}
\usepackage{needspace}
\newtcolorbox{remarkbox}[1][Note]{colback=gray!5,colframe=gray!60,
 coltitle=black,title=#1,fonttitle=\bfseries,arc=2mm,boxrule=0.6pt,breakable}
\newtcolorbox{todobox}[1]{colback=yellow!8,colframe=orange!75!black,
 title={TODO: #1},fonttitle=\bfseries,arc=2mm,boxrule=0.6pt,breakable}
\newtcolorbox{correctionbox}[1]{colback=red!3,colframe=red!55!black,
 title={Correction: #1},fonttitle=\bfseries,arc=2mm,boxrule=0.6pt,breakable}
\newtcolorbox{conceptbox}[1]{colback=blue!3,colframe=blue!55!black,
 title={Supplement: #1},fonttitle=\bfseries,arc=2mm,boxrule=0.6pt,breakable}
\setlength{\parindent}{0pt}
\setlength{\parskip}{6pt}
\allowdisplaybreaks
\newcommand{\E}{\mathbb E}
\newcommand{\Prob}{\mathbb P}
\newcommand{\R}{\mathbb R}
\DeclareMathOperator{\Var}{Var}
\DeclareMathOperator*{\argmin}{arg\,min}
\title{STAT 581: M-estimation and Maximum Likelihood\\\large Notes 4: Handwritten MLE Derivations and Lecture Complements}
\author{Liang-Cheng Chen}
\date{October 28 handwriting; updated October 30, 2026}
\newtcolorbox{gapbox}[1]{colback=cyan!4,colframe=cyan!50!black,title={Missing from original typed notes: #1},fonttitle=\bfseries,arc=2mm,boxrule=.6pt,breakable}
\begin{document}
\maketitle
\begin{remarkbox}[Sources and editorial conventions]
This volume integrates all five substantive pages of the October 28
handwritten \texttt{../notes/stat581-note4.pdf}; page 6 is blank.
Sections on likelihood, consistency, smooth asymptotics, and sandwich
inference incorporate that handwriting. Later sections retain the
previously prepared complements from Yen-Chi Chen's \emph{Lecture 3:
M-estimation and the maximum likelihood estimator}, pp. 3-1--3-26,
including Appendices A--B. They extend beyond this five-page handwritten source.

Gray boxes explain marginal questions, red boxes show corrections, and
blue boxes supply additional arguments. No explicit TODO marker occurs
in this new handwritten file. Handwritten M1--M3 labels below are local
labels and differ from the numbering in lecture Theorem 3.1.
All mathematics uses equation or align environments.
\end{remarkbox}
\clearpage
\tableofcontents
\clearpage
\section{Likelihood, population targets, and misspecification}
\textit{Handwritten p. 1; Lecture 3, Sections 3.1--3.2, pp. 3-1--3-7.}
For IID data under their true law, define the working model and criterion:
\begin{align}
 L_n(\theta)&=\prod_{i=1}^np_\theta(X_i),\quad
 \ell_\theta(x)=\log p_\theta(x),\\
 M_n(\theta)&=\frac1n\sum_i\ell_\theta(X_i),\quad
 M(\theta)=\E_P\ell_\theta(X),\\
 \widehat\theta_n&\in\operatorname*{argmax}_\theta M_n(\theta),\quad
 \theta^*\in\operatorname*{argmax}_\theta M(\theta).
\end{align}
The expectation is under the true law, even if the model is incorrect.
When the relevant KL divergences are defined, maximizing expected
log-likelihood minimizes the KL divergence from the true law to the
working model. With a correct identifiable model the target is the true
parameter. Pointwise LLN alone does not prove convergence of maximizers.
An interior differentiable maximizer satisfies the score equation; a
boundary maximizer or an unbounded likelihood need not do so.

Empirical risk minimization uses the negative criterion. A Z-estimator
instead approximately solves an estimating equation. A stationary point
need not be a global maximizer.
\begin{align}
 R_n(\theta)&=\frac1n\sum_iL(\theta,X_i),\quad R(\theta)=\E L(\theta,X),\\
 S_n(\theta)&=\frac1n\sum_i s_\theta(X_i),\quad
 s_\theta(x)=\nabla_\theta\ell_\theta(x),\quad S(\theta)=\E s_\theta(X),\\
 H_n(\theta)&=\frac1n\sum_i\nabla_\theta^2\ell_\theta(X_i),\quad
 H(\theta)=\E\nabla_\theta^2\ell_\theta(X).
\end{align}
Interchanging derivatives and expectation requires its own domination or
integrability justification.


\begin{remarkbox}[Randomness and the population MLE (handwritten pp. 1--2)]
The working density and the true data law need not coincide. The population
criterion and its unique maximizer are deterministic functions of the true
law and the chosen model; the empirical criterion and its maximizer depend
on the data. Thus the pseudo-true target may change when the working model
changes, even though the true law stays fixed.
\end{remarkbox}
\begin{correctionbox}{A pointwise SLLN is not a uniform law (handwritten p. 1)}
For each fixed parameter, integrability gives the usual SLLN. Its quantifier
does not allow the parameter to vary with the sample or justify exchanging
limit and maximization. A uniform law for the log-likelihood class, together
with separation and approximate maximization, supplies the missing argument.
A standard sufficient setting is compact parameter space, almost-sure
continuity in the parameter, and an integrable envelope for the absolute
log-likelihood. The classical empirical-CDF Glivenko--Cantelli theorem
alone is not a uniform law for an arbitrary likelihood class.
\end{correctionbox}
\begin{remarkbox}[Two meanings of score (handwritten p. 2)]
Here the score differentiates the log-likelihood with respect to the model
parameter. In some machine-learning contexts, ``score'' instead means a
gradient with respect to the data. These are different objects:
\begin{equation}
 s_\theta(x)=\nabla_\theta\log p_\theta(x),\qquad
 s^{\mathrm{data}}_\theta(x)=\nabla_x\log p_\theta(x).
\end{equation}
The latter requires a differentiable continuous-data model. An interior
smooth optimum satisfies a score equation, but a score root alone need
not be the desired maximizer.
\end{remarkbox}

\section{Consistency before normality}
\textit{Theorems 3.7--3.8, pp. 3-8--3-9.}
A useful argmax theorem assumes uniform convergence, a separated maximum,
and sufficiently accurate empirical optimization:
\begin{align}
 &\sup_{\theta\in\Theta}|M_n(\theta)-M(\theta)|=o_P(1),\\
 &\forall\varepsilon>0:\quad
 \sup_{\|\theta-\theta^*\|\ge\varepsilon}M(\theta)<M(\theta^*),\\
 &M_n(\widetilde\theta_n)\ge M_n(\theta^*)-o_P(1).
\end{align}
Then the approximate maximizer is consistent. To prove it, call the
strict population gap outside the target neighborhood a positive number.
Outside that neighborhood, two uniform errors and the optimization error
would have to exceed the gap. The probability of this tends to zero.
Compactness, continuity, and a unique maximum are common sufficient
conditions for separation; uniqueness alone on an arbitrary noncompact
space is insufficient.

The Z-estimator analogue replaces maximum separation by a norm bound:
\begin{align}
 &\sup_\theta\|S_n(\theta)-S(\theta)\|=o_P(1),\quad S(\theta^*)=0,\\
 &\inf_{\|\theta-\theta^*\|\ge\varepsilon}\|S(\theta)\|>0
 \quad\text{for every }\varepsilon>0,\\
 &\|S_n(\widetilde\theta_n)\|=o_P(1).
\end{align}
Uniform convergence and the approximate-zero condition preclude a root
far from the unique separated population zero. In vector problems the
norms are essential; a vector cannot be ordered by a scalar greater-than sign.

\section{Classical smooth asymptotic normality}
\textit{Theorem 3.1 and Remark 3.2, pp. 3-3--3-5.}
Use an interior, identifiable target; consistency from the preceding
argmax theorem; a score equation; an invertible limiting Hessian; a finite
score second moment; and local uniform Hessian convergence. An integrable
absolute envelope on third derivatives is one convenient way to control
Hessian changes. Assume the log-likelihood and its lower derivatives have
the needed integrability at an anchor point.
\begin{correctionbox}{The listed M1--M4 need supporting conditions}
Compactness and a third-derivative envelope alone do not imply finite score
variance or integrability of lower derivatives. A random linear term has
zero third derivative but may have infinite variance. The proof also needs
consistency before its Taylor remainder is negligible. We state these
requirements explicitly rather than infer them from M1--M4 alone.
\end{correctionbox}
Set the Hessian and score covariance at the target as
\begin{equation}
 H=H(\theta^*),\quad B=\E[s_{\theta^*}(X)s_{\theta^*}(X)^{\mathsf T}],
 \quad\E s_{\theta^*}(X)=0.
\end{equation}
A multivariate-safe expansion integrates the Hessian along a segment:
\begin{align}
 0&=S_n(\theta^*)+A_n(\widehat\theta_n-\theta^*),\\
 A_n&=\int_0^1H_n(\theta^*+t(\widehat\theta_n-\theta^*))\,dt
 \xrightarrow{P}H,\\
 \sqrt n(\widehat\theta_n-\theta^*)
 &=-H^{-1}\frac1{\sqrt n}\sum_i s_{\theta^*}(X_i)+o_P(1)\\
 &\xrightarrow{d}N_d(0,H^{-1}BH^{-\mathsf T}).
\end{align}
The local segment is in the parameter space with probability tending to
one because the target is interior and the estimator is consistent.
This expansion avoids assuming a single intermediate point works for an
entire vector score. For symmetric Hessians the rightmost inverse transpose
is just the inverse.
\begin{correctionbox}{Sign and the Taylor remainder}
The score-based influence function is the negative inverse Hessian times
the score. Theorem 3.1's final expansion and Theorem 3.9 omit this minus
sign. It cancels from covariance but matters in linear representations.
A second-order score remainder uses the integral weight corresponding to
Taylor's theorem; a schematic unweighted integral is only a bound, not an
exact remainder formula. The integrated-Hessian identity above is exact.
\end{correctionbox}


\subsection{The handwritten proof, step by step}
\textit{Handwritten pp. 2--5; lecture Theorem 3.1.}
The handwritten assumptions are: M1, both empirical and population targets
solve their score equations; M2, the population target is interior; M3,
entrywise Hessian integrability at the target and nonsingularity of its
expectation. To complete the proof, also assume consistency, finite score
second moment, and local uniform control of the Hessian, as stated above.
These local handwritten labels are not the lecture's M1--M4 list.

\textbf{Step 1: fix the target before applying the CLT.}
Define the empirical fluctuation at the deterministic population target:
\begin{align}
 Y_i&=s_{\theta^*}(X_i),\qquad \E Y_i=0,\\
 \Delta_n&=S_n(\theta^*)-S(\theta^*)
 =\frac1n\sum_{i=1}^n(Y_i-\E Y_i),\\
 B&=\operatorname{Cov}(Y_i)=\E Y_iY_i^{\mathsf T},\qquad
 \sqrt n\Delta_n\xrightarrow{d}N_d(0,B).
\end{align}
At the fixed target these are IID vectors. Evaluating every summand at the
same data-dependent estimator generally makes them dependent, so a direct
IID CLT cannot be applied to those plug-in terms. The empirical score at
the fitted estimator is in fact zero under M1.

\textbf{Step 2: use the score equations and Taylor expansion.}
Write the estimation error and expand about the population target:
\begin{align}
 h_n&=\widehat\theta_n-\theta^*,\\
 -\Delta_n&=S_n(\widehat\theta_n)-S_n(\theta^*)
 =H_n(\theta^*)h_n+R_n,\\
 h_n&=-H_n(\theta^*)^{-1}\Delta_n
       -H_n(\theta^*)^{-1}R_n.
\end{align}
\begin{correctionbox}{Track the remainder sign (handwritten pp. 3--4)}
If the remainder is added in the preceding Taylor expansion, it is
subtracted after solving for the estimation error. Its sign may be
redefined, but both displayed equalities must then use that convention.
It cannot simply be set to zero. The integrated-Hessian argument above
justifies its negligibility and works for a vector score.
\end{correctionbox}
For example, an integrable local envelope for absolute third derivatives
bounds the remainder by a bounded-in-probability multiple of the squared
error norm. Consistency first makes this remainder small relative to the
error norm; the score CLT and nonsingular curvature then yield the
root-sample-size rate. Only after that step is the remainder negligible
at the CLT scale:
\begin{align}
 \|R_n\|&\le C_n\|h_n\|^2,\qquad C_n=O_P(1),\\
 h_n=o_P(1)&\ \Longrightarrow\ R_n=o_P(\|h_n\|),\\
 \|h_n\|&=O_P(n^{-1/2}),\qquad \sqrt n R_n=o_P(1).
\end{align}
The rate follows by absorbing the small multiple of the error norm into
the left side of the solved Taylor inequality. It is not an assumption
that can be silently used to prove itself.

\textbf{Step 3: average IID Hessian matrices and invert the limit.}
For every coordinate pair assume the indicated absolute integrability:
\begin{align}
 W_i&=\nabla_\theta^2\ell_{\theta^*}(X_i),\qquad
 \E |(W_i)_{jk}|<\infty,\\
 H_n(\theta^*)&=\frac1n\sum_iW_i\xrightarrow{P}\E W_i=H,\\
 \det H&\ne0\quad\Longrightarrow\quad
 H_n(\theta^*)^{-1}\xrightarrow{P}H^{-1}.
\end{align}
The last assertion holds on an event with probability tending to one.
The absolute value belongs inside the expectation in the LLN assumption.
In fixed dimension, nonsingularity means every eigenvalue is nonzero;
for a smooth interior maximum the symmetric Hessian is negative
semidefinite, and nonsingularity then makes it negative definite.

\textbf{Step 4: apply Slutsky's theorem.}
The resulting covariance and influence function are
\begin{align}
 \sqrt n h_n&=-H^{-1}\sqrt n\Delta_n+o_P(1)
 \xrightarrow{d}N_d(0,\Sigma),\\
 \Sigma&=H^{-1}BH^{-\mathsf T},\qquad
 \operatorname{IF}(x)=-H^{-1}s_{\theta^*}(x).
\end{align}
The inverse transpose is required by the general covariance transformation;
for a symmetric Hessian it equals the ordinary inverse.

\subsection{Sandwich inference and correct specification}
A consistent plug-in covariance for the scaled estimator is
\begin{align}
 \widehat H&=H_n(\widehat\theta_n),\quad
 \widehat B=\frac1n\sum_i s_{\widehat\theta_n}(X_i)
 s_{\widehat\theta_n}(X_i)^{\mathsf T},\\
 \widehat\Sigma&=\widehat H^{-1}\widehat B\widehat H^{-\mathsf T},\\
 \widehat\theta_{n,j}&\ \pm\ z_{1-\alpha/2}\sqrt{\widehat\Sigma_{jj}/n}.
\end{align}
Consistency of these plug-ins needs suitable local uniform laws and
continuity, not merely convergence of the estimator. Under a correct
regular model, the information identity yields
\begin{equation}
 B=I(\theta_0)=-H,\qquad\Sigma=I(\theta_0)^{-1}.
\end{equation}
Under misspecification, the outer-product matrix is the score variability
under the true law and generally differs from negative expected curvature.
Bootstrap replicate covariance estimates the covariance of the estimator
itself; multiply it by the sample size to compare it with the asymptotic
covariance of the root-sample-size scaled estimator. Bootstrap validity
also requires regularity; it is not universal for every M-estimator.


\begin{correctionbox}{Information identity has a minus sign (handwritten p. 5)}
Under correct specification and differentiation-under-the-integral
conditions, differentiation of the unit total mass gives
\begin{align}
 0&=\nabla_\theta^2\int p_\theta(x)\,dx\\
 &=\E_\theta\!\left[
 \nabla_\theta^2\ell_\theta(X)+s_\theta(X)s_\theta(X)^{\mathsf T}
 \right],\\
 I(\theta_0)&=\E_{\theta_0}s_{\theta_0}(X)s_{\theta_0}(X)^{\mathsf T}
 =-\E_{\theta_0}\nabla_\theta^2\ell_{\theta_0}(X)=-H,\\
 \Sigma&=I(\theta_0)^{-1}=-H^{-1}.
\end{align}
The positive score outer-product expectation is not the expected
log-likelihood Hessian. The handwriting misses this sign in its
correct-model remark. With misspecification, the outer product under the
true law is a score-variability matrix; calling it Fisher information must
not be taken to imply the information identity.
\end{correctionbox}
\begin{remarkbox}[Which covariance does the sandwich estimate? (handwritten p. 5)]
The displayed sandwich estimates the covariance of the limiting scaled
estimator. Divide it by sample size for an estimated covariance of the
unscaled fitted parameter. Plug-in Hessians and outer products require
local uniform convergence, not just a pointwise LLN at a fixed parameter.
In a misspecified model the sandwich is needed even when an inverse-Hessian
formula is available computationally.
\end{remarkbox}

\section{Examples of empirical risk minimization}
\textit{Examples 3.3--3.6, pp. 3-6--3-7.}
Least squares minimizes average squared residuals. For linear prediction,
the population target is a best linear projection, even if the conditional
mean is nonlinear:
\begin{equation}
 \theta^*=\operatorname*{argmin}_\theta\E(Y-X^{\mathsf T}\theta)^2.
\end{equation}
Logistic regression uses the loss and probability
\begin{align}
 q_\theta(x)&=(1+e^{-x^{\mathsf T}\theta})^{-1},\\
 L(\theta;x,y)&=-yx^{\mathsf T}\theta+\log(1+e^{x^{\mathsf T}\theta}),\\
 \nabla L&=x(q_\theta(x)-y),\quad
 \nabla^2 L=xx^{\mathsf T}q_\theta(x)(1-q_\theta(x)).
\end{align}
A finite unique minimizer is an assumption; complete separation can destroy
existence of an unrestricted logistic MLE.

Classification with zero--one loss is also ERM, but the criterion is usually
nondifferentiable, so ERM membership alone does not imply the classical
smooth normality theorem applies.

For a kernel density estimate and its mode,
\begin{align}
 \widehat p_h(x)&=\frac1{nh^d}\sum_iK((X_i-x)/h),\\
 \widehat m_h&\in\operatorname*{argmax}_x\widehat p_h(x),\quad
 m_h\in\operatorname*{argmax}_x\E\widehat p_h(x).
\end{align}
For fixed bandwidth this is an M-estimator for the smoothed mode. For a
shrinking bandwidth it is a triangular problem: under additional kernel,
density, curvature, and bandwidth conditions the scale is
\begin{equation}
 \sqrt{nh^{d+2}}(\widehat m_h-m_h)\xrightarrow{d}N(0,\Sigma_h^{\mathrm{lim}}).
\end{equation}
The lecture's statement is a schematic rate. Centering instead at the true
mode also requires controlling smoothing bias, often of order squared
bandwidth under appropriate smoothness.

\section{Modern M-estimation: nonsmooth empirical criteria}
\textit{Theorem 3.9, Examples 3.10--3.11, pp. 3-9--3-11.}
A useful normality theorem requires: differentiability at the target for
almost every observation; a local Lipschitz envelope with finite second
moment; a nonsingular second-order population expansion; consistency; and
an empirical criterion within a term smaller than inverse sample size of
the empirical optimum (exact maximizers satisfy this).
\begin{align}
 M(\theta)&=M(\theta^*)+\tfrac12(\theta-\theta^*)^{\mathsf T}
 H(\theta-\theta^*)+o(\|\theta-\theta^*\|^2),\\
 |\ell_{\theta_1}(x)-\ell_{\theta_2}(x)|&\le\eta(x)\|\theta_1-\theta_2\|,
 \quad\E\eta(X)^2<\infty.
\end{align}
Then the same negative-Hessian score expansion applies. The sample
criterion need not have a second derivative everywhere.

\subsection{Sample median}
Assume a unique population median with a density continuous and positive
there. For a finite expected absolute loss, additionally assume a finite
first moment (alternatively work with centered loss differences).
\begin{align}
 \ell_\theta(x)&=-|x-\theta|,\quad m=F^{-1}(1/2),\\
 s_m(x)&=\operatorname{sign}(x-m)\quad(x\ne m),\\
 M'(\theta)&=1-2F(\theta),\quad H=M''(m)=-2f(m),\\
 \psi(x)&=\frac{\operatorname{sign}(x-m)}{2f(m)},\\
 \sqrt n(\widehat m-m)&\xrightarrow{d}
 N\left(0,\frac1{4f(m)^2}\right).
\end{align}
The local Lipschitz envelope is one. An observation equal to the target
has probability zero, so its kink does not violate almost-everywhere
differentiability at that target. A consistent density plug-in or a valid
median bootstrap supplies uncertainty estimates.
\begin{correctionbox}{Median maximization has negative curvature}
Example 3.10 reverses the signs in the absolute-loss integrals, score,
and Hessian. For negative absolute loss the population maximum has
negative curvature. The displayed limiting variance in the lecture is
nevertheless correct because the signs cancel when squared.
\end{correctionbox}

\subsection{Nonlinear least squares}
Suppose the conditional mean model is correct, errors have conditional
mean zero, and the needed differentiability and integrability hold:
\begin{align}
 Y&=m_{\theta^*}(X)+e,\quad \E[e\mid X]=0,\quad
 g(X)=\nabla_\theta m_{\theta^*}(X),\\
 \ell_\theta(X,Y)&=-\tfrac12(Y-m_\theta(X))^2,\\
 s_{\theta^*}(X,Y)&=e\,g(X),\quad
 A=\E[g(X)g(X)^{\mathsf T}],\quad H=-A,\\
 B&=\E[e^2g(X)g(X)^{\mathsf T}],\quad
 \Sigma=A^{-1}BA^{-1}.
\end{align}
Under constant conditional variance this reduces to sampling variance
times the inverse of the Gram matrix. The half in the criterion ensures
the displayed derivative has no factor two. Without it both the score
and Hessian double, leaving the sandwich unchanged.
\begin{correctionbox}{Curvature is not automatically Fisher information}
Example 3.11 drops derivative factors and the negative curvature sign.
The Gram matrix is not generally Fisher information for an arbitrary
error law. A likelihood and error distribution are required to make that
identification. The constant population criterion at the truth is negative
one-half the error second moment with the convention used here.
\end{correctionbox}

\section{Influence functions and deletion diagnostics}
\textit{Section 3.3.3, p. 3-12.}
For the maximization convention,
\begin{equation}
 \widehat\theta_n-\theta^*=\frac1n\sum_i\psi(X_i)+o_P(n^{-1/2}),
 \quad\psi(x)=-H^{-1}s_{\theta^*}(x).
\end{equation}
Large influence indicates sensitivity. For the corresponding linearized
estimators, deleting one observation changes the average by
\begin{equation}
 \widehat\theta_n-\widehat\theta_{n,-k}
 \approx\frac1n\left\{\psi(X_k)-\frac1{n-1}\sum_{i\ne k}\psi(X_i)\right\}.
\end{equation}
\begin{correctionbox}{The deletion effect has an inverse-sample-size factor}
Section 3.3.3 omits this factor. The influence function describes a scaled
first-order sensitivity. The ordinary asymptotic-linear remainder alone
is too coarse to prove an inverse-sample-size deletion approximation;
this interpretation requires additional smoothness or a direct deletion
analysis. The displayed identity is exact for the linearized averages.
\end{correctionbox}

\section{Quadratic mean differentiability and LAN}
\textit{Section 3.4, Theorems 3.13--3.14, pp. 3-12--3-17.}
QMD is differentiability of the square-root density in the square-integrable
function space with respect to a common dominating measure:
\begin{equation}
 \int\left[\sqrt{p_{\theta+h}}-\sqrt{p_\theta}
 -\tfrac12h^{\mathsf T}s_\theta\sqrt{p_\theta}\right]^2d\mu
 =o(\|h\|^2).
\end{equation}
The full-support formulation is essential when support changes. A formal
pointwise derivative on the interior of the support is insufficient.
QMD implies zero-mean square-integrable score under the model. In the
normalized Hellinger convention of Lecture 2,
\begin{equation}
 H^2(P_{\theta+h},P_\theta)=\tfrac18h^{\mathsf T}I(\theta)h
 +o(\|h\|^2).
\end{equation}
If Hellinger squared is defined without the one-half integral factor,
the coefficient is one-quarter instead.

For a uniform law on an interval with unknown right endpoint, a positive
endpoint shift creates a new support interval whose square-root-density
error is already first order in the shift:
\begin{equation}
 \int_\theta^{\theta+h}p_{\theta+h}(x)\,dx
 =\frac{h}{\theta+h},\qquad h>0.
\end{equation}
Hence the model is not QMD at the usual root-sample-size scale. Indeed,
the endpoint MLE is the sample maximum and has the nonnormal limit
\begin{equation}
 n(\theta-\max_iX_i)/\theta\xrightarrow{d}\operatorname{Exp}(1).
\end{equation}
This last limit is an editorial calculation: use the exact CDF of the maximum.

For a correct QMD model, with bounded local vectors tending to a fixed
vector, the LAN expansion under the target model is
\begin{align}
 \log\prod_i\frac{p_{\theta+h_n/\sqrt n}(X_i)}{p_\theta(X_i)}
 &=h_n^{\mathsf T}\Delta_n-\tfrac12h_n^{\mathsf T}I(\theta)h_n+o_P(1),\\
 \Delta_n&=\frac1{\sqrt n}\sum_i s_\theta(X_i)
 \xrightarrow{d}N(0,I(\theta)).
\end{align}
For a consistent MLE with nonsingular information and the local
log-likelihood Lipschitz condition of Theorem 3.14 (with an envelope having
a finite second moment),
\begin{equation}
 \sqrt n(\widehat\theta_n-\theta_0)
 =I(\theta_0)^{-1}\Delta_n+o_P(1)
 \xrightarrow{d}N(0,I(\theta_0)^{-1}).
\end{equation}
\begin{correctionbox}{Hellinger, KL, and proof qualifications}
The source switches Hellinger normalizations. Fixing the convention avoids
a factor-two error. Its equation (3.11) equates the expected log-likelihood
difference with minus twice an unnormalized Hellinger square; that is only
a local second-order relation when an expected log-likelihood expansion
is justified, not an exact identity. QMD controls square-root densities;
passing a pointwise logarithm remainder through expectation requires more
control. Use the LAN theorem and the full assumptions of Theorem 3.14,
rather than treating that informal proof as valid from QMD alone.
\end{correctionbox}

\section{Many nuisance parameters: the missing-outcome example}
\textit{Section 3.5, pp. 3-17--3-18.}
Let there be many categories, with known observation probabilities bounded
away from zero and one:
\begin{align}
 X&\sim\operatorname{Unif}\{1,\ldots,B\},\quad
 R\mid X=b\sim\operatorname{Bernoulli}(\xi_b),\\
 Y\mid X=b&\sim\operatorname{Bernoulli}(\theta_b),\quad
 R\perp Y\mid X,\quad\delta\le\xi_b\le1-\delta,\\
 \psi_B&=B^{-1}\sum_{b=1}^B\theta_b.
\end{align}
Only outcomes with indicator one are observed. In categories with observed
outcomes, the MLE is the empirical success proportion; in empty categories
the likelihood is flat. With far more categories than observations, most
coordinates are unidentified by the observed likelihood and an arbitrary
plug-in completion or independent prior can dominate the result.

An inverse-probability estimator targets the functional directly:
\begin{align}
 Z_i&=R_iY_i/\xi_{X_i},\quad\widehat\psi_B=n^{-1}\sum_iZ_i,\\
 \E Z_i&=B^{-1}\sum_b\theta_b=\psi_B,\quad
 \E Z_i^2=B^{-1}\sum_b\theta_b/\xi_b\le1/\delta,\\
 \Var(\widehat\psi_B)&\le1/(n\delta).
\end{align}
Chebyshev proves consistency even along growing-category sequences.
A triangular-array CLT also gives asymptotic normality if the row variances
converge to a positive limit; bounded summands verify Lindeberg. Without
variance convergence there need not be a single fixed normal limit.
\begin{correctionbox}{Interpret the failure claim in its asymptotic regime}
With fixed finitely many categories and positive observation probabilities,
every category is eventually observed and standard consistency can return.
The difficulty is the growing or very large category regime, not a universal
failure of every Bayesian prior or penalized estimator. A penalized MLE is
a posterior mode for a matching penalty, not generally the posterior-mean
Bayes rule. Estimating the average can succeed without recovering every
coordinate.
\end{correctionbox}

\section{Computational learning: gradient descent}
\textit{Appendix A, Theorem 3.17, pp. 3-19--3-21.}
On a convex region suppose a differentiable objective has a Lipschitz
gradient and a positive strong-convexity constant:
\begin{align}
 f(y)&\le f(x)+\nabla f(x)^{\mathsf T}(y-x)+\tfrac L2\|y-x\|^2,\\
 f(y)&\ge f(x)+\nabla f(x)^{\mathsf T}(y-x)+\tfrac M2\|y-x\|^2.
\end{align}
When twice differentiable, these correspond to all Hessian eigenvalues
between the two constants. Gradient descent updates
\begin{equation}
 \theta^{(t+1)}=\theta^{(t)}-\gamma\nabla f(\theta^{(t)}).
\end{equation}
For a stationary minimizer and a fixed step smaller than both reciprocal
constants, the proof in the lecture yields
\begin{equation}
 \|\theta^{(t)}-\widehat\theta_n\|^2
 \le(1-\gamma M)^t\|\theta^{(0)}-\widehat\theta_n\|^2.
\end{equation}
The argument expands the squared update, uses strong convexity to control
the cross term, and smoothness to bound squared gradient by twice the
smoothness constant times excess objective value. The remaining gradient
term is nonpositive for the stated step size.

\section{Transfer population curvature to the empirical objective}
\textit{Appendix B, Lemmas 3.18--3.19 and Theorem 3.20, pp. 3-22--3-26.}
For the negative population criterion, suppose the target Hessian has
smallest eigenvalue bounded below by a positive constant and is locally
Lipschitz in operator norm:
\begin{equation}
 \lambda_{\min}(\nabla^2R(\theta^*))=\lambda_*>0,\quad
 \|\nabla^2R(\theta)-\nabla^2R(\theta^*)\|\le\phi\|\theta-\theta^*\|.
\end{equation}
Choose a ball lying in the interior of the parameter space, of radius at
most the smallest eigenvalue divided by twice the Lipschitz constant.
Weyl's eigenvalue inequality then ensures curvature at least half the
target curvature throughout that ball.

A compact Lipschitz-indexed class has a uniform LLN when its envelope of
increments is integrable and one anchor function is integrable:
\begin{equation}
 |f_{\theta_1}(x)-f_{\theta_2}(x)|\le q(x)\|\theta_1-\theta_2\|,
 \quad\E q(X)<\infty,\quad\E|f_{\theta_a}(X)|<\infty.
\end{equation}
A finite parameter net, ordinary LLN at its finitely many points, and a
Lipschitz approximation bound prove this. Applying it to Hessian entries
in fixed dimension yields uniform convergence in operator norm.

On an event with probability tending to one, require all of the following:
the empirical estimator is within half the population ball radius of the
target; the uniform Hessian error is at most one-quarter the target
curvature; and empirical Hessian norms are bounded by twice a finite
population bound. Then in a ball around the empirical estimator,
\begin{equation}
 M^*=\lambda_*/4,\qquad L^*=2h_{\max},\qquad
 \|\theta^{(t)}-\widehat\theta_n\|
 \le(1-\gamma M^*)^{t/2}\|\theta^{(0)}-\widehat\theta_n\|.
\end{equation}
The contraction keeps iterates in the local ball if initialization lies
there. There is no guarantee that an arbitrary initial point reaches this
ball. The lecture's last proof uses the maximum-eigenvalue subscript in
several lower bounds where the minimum-eigenvalue subscript is needed;
the union bound must count each of its three events once.

\subsection{Computational error versus statistical error}
The two errors add by the triangle inequality:
\begin{equation}
 \|\theta^{(t)}-\theta^*\|
 \le\|\theta^{(t)}-\widehat\theta_n\|
 +\|\widehat\theta_n-\theta^*\|.
\end{equation}
A fixed contraction factor requires logarithmically many iterations in
sample size to reach the statistical root-sample-size scale. To preserve
the estimator's limiting distribution, require computational error to be
smaller in probability than that scale, not merely bounded by it. This
last comparison is an editorial synthesis of the two lecture appendices.

\end{document}
